Definieren Sie VWL als Teilbereich

Die Volkswirtschaftslehre wird auch Nationalökonomie oder Sozialökonomie genannt
und ist ein Teilgebiet der Wirtschaftswissenschaften.
Untersuchungsgegenstände der Volkswirtschaftslehre sind jeweils Zusammenhänge
in der Produktion und Allokation von Gütern und Produktionsfaktoren.

Mögliche Entscheidungen die von Gesellschaften zu treffen sind:
- Welche Güter sollen von wem und in welchem Umfang produziert werden?
- Welche Ressourcen sollen für die Produktion eingesetzt werden?
- Wie soll der Wohlstand verteilt werden?

In der volkswirtschaftlichen Theorie unterscheidet man grundsätzlich
zwischen den folgenden Ressourcen: Kapital, Arbeit und Land.
• Diese Ressourcen sind allerdings in der realen Welt begrenzt, sodass
auch die Volkswirtschaftslehre in ihren Modellen von begrenzten
Ressourcen ausgeht. Gegenstand der Volkswirtschaftslehre ist damit
das „managen“ der knappen Ressourcen und finden von „optimalen“
Lösungen.
• Gegenstand ist daher die Lösung des Knappheitsproblems durch eine
effiziente Allokation der vorhandenen Ressourcen.
• Die Effizienz wird anhand einer Zielfunktion bestimmt.


Wie unterscheiden sich die Mikro- und Makroökonomik?

Die Mikroökonomik befasst sich mit der (volkswirtschaftlichen) Analyse des Verhaltens von
Individuen und deren wirtschaftlichen Verhaltens. 
Die betrachteten Subjekte sind dabei die
Haushalte und die Unternehmen. 
Fragestellung der Mikroökonomik ist insbesondere
- wie die Interaktion zwischen den Subjekten (z.B. Angebot von Arbeitsleistung durch die
Haushalte oder Produktion von Gütern durch die Unternehmen) erfolgt
- Wie knappe Ressourcen und Güter allokiert werden 
- Hauptgegenstand ist der Marktmechanismus und die Preisbildung


Die Makroökonomik befasst sich mit der übergreifenden
volkswirtschaftlichen Gesamtbetrachtung und den
entsprechenden Komponenten und Akteuren.

- Wesentliche Komponenten sind bspw. die Geld-, und
Fiskalpolitik
- Wichtigste Teilgebiete sind die volkswirtschaftliche
Gesamtrechnung, Einkommens- und Beschäftigungs-
theorie, Wachstumstheorie und die Konjunkturtheorie


Was sind normative und positive Aussagen der VWL


Die normative Ökonomik befasst sich mit Diskursen über
Wirtschaftsordnung, politischen Zielen und Absichten.
Sie basiert daher auf Wertvorstellungen und sozialen Normen und
ist präskriptiv (verordnend).
So trifft beispielsweise die Wohlfahrtsökonomik
Aussagen darüber, ob bestimmte Veränderungen wünschenswert sind.

Die normative Analyse in der Volkswirtschaftslehre beschäftigt sich
vor allem mit der Frage „Was sollte sein?“. Diese Analyse ist eine
über eine Prognose hinausgehende Bewertung, die z. B. auf die
Frage „Was ist das Beste?“ näher eingeht.
De normative Analyse fragt beispielsweise
„Sollte die Arbeitslosigkeit reduziert werden?“.

Die positive Ökonomik umfasst die Analysen und Theorien über die
ökonomische Welt, behandelt also deskriptiv die Fakten. Es handelt
sich um eine positive Wissenschaft. Die Gültigkeit einer derart
getroffenen Aussage kann beispielsweise durch statistische Daten
widerlegt werden.

Die positive Analyse in der Volkswirtschaftslehre beschäftigt sich
mit den ökonomischen Konsequenzen eines bestimmten Ereignisses
bzw. einer bestimmten politischen Maßnahme. Dies geschieht
unabhängig davon, ob diese Konsequenzen wünschenswert sind
oder nicht. Die Analyse erfolgt unter der Leitfrage „Was ist und
warum ist es so, wie es ist?“. Eine Frage der positiven Analyse ist
demnach „Wie wird sich die Arbeitslosigkeit entwickeln?“.



Nennen Sie die 9 volkswirtschaftlichen Regeln

Regel 1 – Alle Menschen müssen sich zwischen Alternativen entscheiden

Regel 2 - Die Kosten eines Gutes ergeben sich aus den Opportunitätskosten
Der Nutzen ist individuell auf Basis der verschiedenen Präferenzen gebildet und kann anhand einer sogenannten Nutzenfunktion ausgedrückt werden
• Wie werden Kosten einheitlich und vergleichbar bestimmt?
• Als Opportunitätskosten wird das verstanden was es an Nutzen kostet, eine Möglichkeit nicht durchzuführen


Regel 3 – Rationale Menschen treffen Grenzentscheidungen
Die VWL betrachtet insbesondere die Konzepte der Grenzkosten sowie des Grenzertrag
- Grenzkosten: Was kostet es mich eine (marginale) Einheit mehr zu produzieren
- Grenzertrag: Was bringt die eine weitere (marginale) produzierte und verkaufte Einheit an zusätzlichem Ertrag
Wenn gilt: Grenzkosten < Grenzertrag, führt die Produktion (und Verkauf) einer weiteren Einheit eines Gutes zu einem Gewinn. 
Dies ist solange der Fall, bis Grenzkosten = Grenzertrag oder Grenzkosten > Grenzertrag gilt.

Regel 4 –Menschen reagieren auf Anreize

Ein sehr wichtiges Instrument zur Messung von voraussichtlichen Verhaltensänderungen ist 
die Elastizität der Nachfrage bzw. Elastizität des Angebots.
Steigt die Nachfrage bei gegebenen Preisänderungen überproportional, spricht man von einer sehr elastischen Nachfrage.
Ändern Konsumenten und Produzenten ihre Nachfrage bzw. ihr Angebot überhaupt nicht, ist die Nachfrage bzw. das Angebot vollkommen unelastisch.

Sei eine differenzierbare Nachfragefunktion \( x(p) \). Die Preiselastizität der Nachfrage ist definiert:
\( \nu_{x,p} = \frac{\frac{dx}{x}}{\frac{dp}{p}} \)

Im Bereich des Versicherungsschutzes dürfte die Nachfrage nach Absicherung existentieller Risiken 
(Lebensversicherung, Krankenversicherung,…) unelastischer sein als die Nachfrage nach Schutz gegen überschaubare Risiken 
(z.B. Reiseversicherung, Krankenhaustagegeld,…)

Regel 5 – Durch Handel kann es allen Beteiligten besser gehen
Land Tuch Wein
A
2h/Einheit 5h/Einheit
B
4h/Einheit 2h/Einheit


Regel 6 – Märkten gelingt es gewöhnlich gut, das Wirtschaftsleben zu organisieren
In der (konkurrierenden) Marktwirtschaft gibt es keine zentrale Planungsbehörde, 
das Wirtschaftsleben ist das Resultat der (Individual-) Entscheidungen
der Wirtschaftssubjekte, die eigenständig über ihren Konsum und ihre Produktion entscheiden.
Angebot und Nachfrage werden über den Markt und insbesondere über den Preismechanismus zusammengeführt

Regel 7 - Der Eingriff einer Regierung kann die Marktergebnisse manchmal verbessern
Der freie Markt führt in der klassischen Theorie zu einer bestmöglichen Gestaltung des Wirtschaftslebens, 
wie kann eine Regierung die Ergebnisse verbessern?
- Sicherstellung von Eigentum
- Beseitigung von Marktversagen aufgrund externer Effekte (bspw. Verunreinigung von Gewässern durch Produktion einer Fabrik)

Regel 8 – Die Wohlstand eines Landes hängt von seinen Möglichkeiten Güter zu produzieren ab
Die Stundenproduktivität ist eine Maßzahl für die Menge an Gütern und Dienstleistungen, 
die innerhalb einer bestimmten Zeit pro Arbeiter bzw. Arbeitnehmer produziert werden können

Regel 9 – Preise steigen, wenn sich zu viel Geld im Wirtschaftskreislauf befindet

Steigende Preise bezeichnet man ökonomisch als Inflation, fallende Preise als Deflation. 
Eine hohe Inflation führt zu Wohlfahrtsverlusten sowie zu ungewollten
Verteilungseffekten, z.B. zwischen Schuldnern und Gläubigern
Den Zusammenhang zwischen Preisniveau und Geldmenge liefert die Quantitätsgleichung:
\( MV = PY \)
Dabei bezeichnet 
- \( M \) die Geldmenge, 
- \( V \) die Umschlaggeschwindigkeit des Geldes, 
- \( P \) das Preisniveau 
- \( Y \) das reale Bruttoinlandsprodukt


Erklaeren Sie Größenbegriffe und deren Bewertung

Stromgroßen
Stromgrößen lassen sich über einen Zeitraum messen. Beispiele:
• Bruttoinlandsprodukt
• Gewinn- und Verlustrechnung eines Unternehmens

Nominale Bewertung
Bewertet eine Auswahl an Gütern mit den jeweils tagesaktuellen Preisen.
Der Wert des Güterbündels wird somit durch die Inflation beeinflusst.

Bestandsgrößen
Bestandsgrößen beziehen sich auf einen Zeitpunkt. Beispiele:
• Arbeitslosigkeit zu einem bestimmten Stichtag
• Bilanz eines Unternehmens für das abgelaufene Geschäftsjahres

Reale Bewertung
Bewertet eine Auswahl an Gütern mit einem definierten, konstanten
Preisniveau. Die Bewertung ist damit um Preisanstiege bereinigt.

Makroökonomische Kennzahlen
Für die Messung der wirtschaftlichen Leistung – und damit Gegenstand der makroökonomischen Analyse – sind verschiedenste Kennzahlen, zu den
wichtigsten zählen:
• die Arbeitslosenquote
• Die Inflationsrate
• Das Bruttoinlandsprodukt


NICHT SO WICHTIG
Erklären Sie die 3 Berechnungsmethoden für das BIP 

BIP Def.
Marktwert aller Waren und Dienstleistungen, die in einem Land, in
einem bestimmten Zeitraum hergestellt werden und für den
Endverbrauch bestimmt sind.

Berechnung
Grundsätzlich werden 3 Berechnungsmethoden unterschieden:
• Entstehungsrechnung
• Verwendungsrechnung
• Verteilungsrechnung


• Ein Autohersteller: dieser kauft Stahl und produziert hieraus
Autos unter Einsatz von Maschinen und Arbeitskraft.
Für ein verkauftes Auto benötigt er eine Tonne Stahl und kann
dieses für 210 EUR verkaufen. Den Stahl kauft er vom
Stahlproduzenten für 100 EUR, Aufwendungen für Löhne
betragen 70 EUR , sodass ein Profit von 40 EUR verbleibt.

BIP nach Entstehungsrechnung:
Die Summe der Wertschöpfung entspricht dem Wert des verkauften Autos sowie des
produzierten Stahls, aber abzüglich der eingesetzten Vorprodukte (Stahl beim Autohersteller):
210 EUR+100 EUR–100 EUR=210 EUR

BIP nach Verwendungsrechnung:
In unserem Beispiel kaufen:
• Private Konsumenten Autos für 210 EUR
• Führen die Unternehmen keine Investitionen in neue Maschinen o.Ä. durch
• Gibt es keine Staatsausgaben und keinen internationalen Warenverkehr
Damit beläuft sich das BIP nach dieser Berechnung ebenfalls auf 210 EUR

BIP nach Verteilungsrechnung:
Das BIP nach der Verteilungsrechnung entspricht der Summe der Einkommen:
80 EUR (Lohn der Stahlarbeit)+70 EUR (Lohn der Autobauer)+60 EUR (Unternehmensgewinne)
=210 EUR
